\section{Why Existing Colocation Abstractions Miss Bubbles}
\label{sec:gap}

\acresetall

Existing inference--fine-tuning colocation systems reclaim capacity
visible at serving scheduling boundaries, such as decode iterations or
periods of sub-peak load. Expert-parallel bubbles are different: they
open and close inside an \ac{MoE} layer and differ across ranks within
the same forward pass. Harvesting them therefore requires not only short
work units, but a planner that knows \emph{which rank} has slack and
\emph{for how long} before launching non-preemptible GPU work.

\subsection{Existing planners reason at serving boundaries}
\label{sec:gap-boundary}

Table~\ref{tab:gap} summarizes the mismatch. Prior systems make admission
decisions from serving-level state such as batch composition, iteration
latency, or engine headroom. Such signals are appropriate for the
capacity they target, but aggregate away the rank-local structure of an
\ac{EP} bubble. Even a perfect prediction of the enclosing iteration's
duration does not reveal which rank will finish expert computation first
or how much earlier it will finish.

\begin{table*}[t]
  \centering
  \small
  \caption{
    \textbf{Bubble harvesting requires finer planning granularity, not only
    finer work units.}
    Prior colocation systems reason at serving scheduling boundaries;
    \sys{} predicts slack per rank and \ac{MoE} layer from the current
    router assignment.
  }
  \label{tab:gap}
  \begin{tabular}{lllll}
    \toprule
    System & Capacity targeted & Work unit & Planning granularity & Signal \\
    \midrule
    FlexLLM~\cite{flexllm}
    & decode underutilization
    & token slot
    & serving iteration
    & batch composition \\
    LLMStation~\cite{llmstation}
    & decode headroom
    & tasklet
    & serving iteration
    & iteration timing \\
    Harli~\cite{harli}
    & decode underutilization
    & SM partition
    & serving iteration
    & predicted latency \\
    DeltaServe~\cite{deltaserve}
    & serving headroom
    & LoRA batch
    & serving iteration
    & engine state \\
    \midrule
    \sys{}
    & \ac{EP} bubble
    & backward chunk
    & rank $\times$ \ac{MoE} layer
    & router top-$k$ assignment \\
    \bottomrule
  \end{tabular}
\end{table*}

The needed signal becomes available immediately after routing: the
serving engine knows the expert work assigned to each rank before expert
dispatch begins. Section~\ref{sec:char-predict} shows that this state
predicts rank-local slack, so \sys{} places its admission decision between
routing and expert execution rather than around the enclosing serving
iteration.

\subsection{Finer work units do not fix the planning mismatch}
\label{sec:gap-granularity}

Smaller fine-tuning chunks help work fit inside millisecond-scale bubbles,
but do not tell the planner where those bubbles occur. We isolate these
two effects with a trace-driven experiment on Qwen3-30B-A3B,
\ac{TP}$=2$/\ac{EP}$=2$, at 128k tokens. We replay the same 486
measured C100 backward chunks (median service time 567\,$\mu$s) against
the same serving trace for three planners. A \emph{router-blind} planner
assigns each layer one historical, rank-independent budget. An
\emph{oracle rank-agnostic} planner receives the true mean bubble across
ranks for the current request and layer, but not which rank owns the
slack. \sys{} instead predicts a separate budget for each rank from the
current router assignment.

We measure \emph{bubble utilization}, the fraction of available bubble
GPU-time occupied by admitted fine-tuning work, and compare each planner
near a 1\% overrun rate.

\begin{table}[t]
  \centering
  \small
  \caption{
    \textbf{Router state closes the planning gap.}
    Using identical measured chunks, router-aware planning recovers
    42.7--45.1\% of available bubble capacity near a 1\% overrun rate,
    versus approximately 1.3\% for rank-agnostic planners.
  }
  \label{tab:gap-planner}
  \begin{tabular}{lrr}
    \toprule
    Planner & Bubble utilization & Overrun rate \\
    \midrule
    Router-blind
    & 1.28--1.31\%
    & 1.04\% \\
    Oracle rank-agnostic
    & 1.22--1.30\%
    & 1.04\% \\
    Router-aware (\sys{})
    & 42.74--45.14\%
    & 0.83--1.04\% \\
    \bottomrule
  \end{tabular}
\end{table}

At comparable overrun frequency, router-aware planning improves harvested
capacity by 33--35$\times$. The oracle performs no better than the
router-blind planner despite knowing the true mean slack for the current
request and layer: knowing \emph{how much} slack exists in aggregate does
not reveal \emph{which rank} owns it.

Execution granularity is a separate constraint. With the coarser C1000
partition, 97 measured chunks have a median service time of 4.13\,ms and
router-aware utilization is 37.95--38.46\%. Recompiling to C100 reduces
the median to 567\,$\mu$s and raises utilization to 42.74--45.14\%,
while the router-blind planner remains near 1--1.5\% under both
partitions. Even an idealized zero-overhead 4$\times$ subdivision of
C1000 chunks increases router-blind utilization only from 1.50\% to
2.09\%, compared with 38.46\% to 49.88\% for router-aware planning.

Thus, shorter chunks determine \emph{how much} work can fit, while router
state determines \emph{where} it can safely execute. \sys{} requires
both, but the latter is the missing colocation abstraction: rank-specific
admission inside the \ac{MoE} forward path.
